\section{From calibrated Gaussian moments to cumulants}
\label{sec:allarity-cumulant-comparison}

The marked comparison uses all subsets of a finite insertion word. Its
Gaussian calculation must therefore control products before taking the
finite logarithm. We give the degree-dependent estimates with the powers
used in the scale purchases. The covariance and calibration hypotheses
are stated explicitly; their common-atlas construction is needed when
this calculation is applied to the selected history.

\subsection{A dimension-independent polynomial estimate}

For a polynomial $p$ in finitely many commuting variables, write
$|p|_1$ for the sum of the absolute values of its coefficients. It is
submultiplicative: expand a product, apply the triangle inequality to
each coefficient, and then sum over the two original monomial indices.
Let $\mathcal W_{Q,\mu}$ be the Wick evaluation with symmetric
covariance $Q$ and mean $\mu$. No positivity is needed in this algebraic
calculation. The defining recurrence, for a product of $j$ coordinate
variables, is
\begin{equation}\label{eq:allarity-wick-recursion}
 \mathcal W_{Q,\mu}(X_{i_1}\cdots X_{i_j})
 =\mu_{i_1}\mathcal W_{Q,\mu}(X_{i_2}\cdots X_{i_j})
  +\sum_{\ell=2}^j Q_{i_1i_\ell}
       \mathcal W_{Q,\mu}\left(\prod_{a\ne1,\ell}X_{i_a}\right),
\end{equation}
with value one on the empty product. Repeated coordinate indices are
allowed: their positions in this recurrence remain distinct.

Assume $B\ge1$, $|\mu_i|\le B$ and $|Q_{ij}|\le B$ for every
coordinate used. Define $C_d(B)=[2B(d+1)]^d$, including $C_0(B)=1$.
Then every polynomial of degree at most $d$ satisfies
\begin{equation}\label{eq:allarity-wick-bound}
           |\mathcal W_{Q,\mu}(p)|\le |p|_1 C_d(B).
\end{equation}
If a second pair $(Q',\mu')$ obeys the same bounds and
$|Q_{ij}-Q'_{ij}|,|\mu_i-\mu'_i|\le\varepsilon$, then
\begin{equation}\label{eq:allarity-wick-difference}
 |\mathcal W_{Q,\mu}(p)-\mathcal W_{Q',\mu'}(p)|
                   \le |p|_1\varepsilon C_d(B).
\end{equation}
Here $\varepsilon\ge0$ can be arbitrary.

We prove both estimates on monomials by simultaneous induction on the
degree ceiling $d$. At degree zero their assertions are respectively
$1\le1$ and $0\le\varepsilon$. For a nonempty monomial of degree at
most $d+1$, the recurrence has one mean term and at most $d$ covariance
terms. Every remaining product has degree at most $d$. Thus the first
bound is at most $(d+1)B C_d(B)\le C_{d+1}(B)$.
For a difference of terms use
$ab-a'b'=a(b-b')+(a-a')b'$. The induction bounds its absolute value
by $B\varepsilon C_d(B)+\varepsilon C_d(B)$, at most
$2B\varepsilon C_d(B)$. Summing at most $d+1$ such terms gives
$2B(d+1)\varepsilon C_d(B)\le\varepsilon C_{d+1}(B)$.
The inequalities follow from
$C_d(B)\le[2B(d+2)]^d$. The empty monomial satisfies both bounds
at every degree ceiling. Finally expand $p$ in monomials and sum their
absolute coefficients. This proves \eqref{eq:allarity-wick-bound}
and \eqref{eq:allarity-wick-difference} without a factor depending on
the number of coordinate variables.

\subsection{Calibrated vertices and their products}

Fix $n$ insertion slots, a common polynomial degree bound $k$, and
possibly different grades $g_i$ for the slots. Here $n,k,g_i$ are
nonnegative integers, and all error parameters $\xi,\varepsilon$ are
nonnegative. After substituting the
actual linear reader map, write $H_i$ for the polynomial of slot $i$.
Let $F_i$ be its intended free polynomial, embedded in the physical
coordinates of the same enlarged variable space. Suppose the grade
truncation of $H_i$ through grade $g_i$ is exactly $F_i$.
Every variable has a nonnegative integer grade. For $r\ge1$, write
\[
 |p|_{1,r}=\sum_\alpha |p_\alpha|r^{\operatorname{grade}(\alpha)}.
\]
Every monomial of $H_i-F_i$ has grade at least $g_i+1$, whence
\begin{equation}\label{eq:allarity-grade-tail}
 r^{g_i+1}|H_i-F_i|_1\le |H_i|_{1,r}.
\end{equation}
This uses the exact low-grade identity, so no difference of two large
low-grade coefficients is being estimated.

Suppose the incoming polynomial has coefficient norm at most $a_i$
and degree at most $k$, and every substituted linear coordinate has
weighted norm at most $b\ge1$. Expanding each monomial and using
submultiplicativity gives $|H_i|_{1,r}\le a_i b^k$. Thus the purchase
\[
                    a_i b^k\le\xi r^{g_i+1}
\]
implies $|H_i-F_i|_1\le\xi$. Also assume $|F_i|_1\le A$, where
$A\ge1$. Then $|H_i|_1\le A+\xi$.
For a subset $I$ of slots, telescope the product one factor at a time.
Every summand contains one difference $H_i-F_i$ and $|I|-1$ factors
whose norm is at most $A+\xi$. Therefore
\begin{equation}\label{eq:allarity-vertex-product}
 \left|\prod_{i\in I}H_i-\prod_{i\in I}F_i\right|_1
 \le |I|\xi(A+\xi)^{|I|-1}
 \le n\xi(A+\xi)^n.
\end{equation}
For $I=\varnothing$ the difference is zero. Polynomial degrees are at
most $nk$, both before and after substitution and truncation.

The actual reader covariance and mean must be the pullbacks of a common
pair $(Q,\mu)$ on the enlarged variable space. Substitution then commutes
with Wick evaluation. To verify this, apply
\eqref{eq:allarity-wick-recursion} to products of substituted linear
forms. Their means are the linear pullbacks of $\mu$ and their pairings
are the bilinear pullbacks of $Q$. Induction on the number of factors
gives the identity on monomials; linearity gives it on all polynomials.
This establishes the needed identity without changing Gaussian variables
separately for each slot.

Assume all entries of $(Q,\mu)$ are bounded by $B\ge1$, and its
restriction to the physical coordinates is within $\varepsilon$ of
the free covariance with zero mean, whose entries are also bounded by
$B$. The $F_i$ involve only these physical coordinates. Insert and
subtract the Wick evaluation of $\prod_{i\in I}F_i$ at $(Q,\mu)$.
Equations \eqref{eq:allarity-wick-bound}--\eqref{eq:allarity-vertex-product}
give, for every subset $I$,
\begin{equation}\label{eq:allarity-moment-error}
 \left|\mathcal W_{Q,\mu}\left(\prod_{i\in I}H_i\right)
       -\mathcal W_{\rm free}\left(\prod_{i\in I}F_i\right)\right|
 \le\{n\xi(A+\xi)^n+A^n\varepsilon\}\,C_{nk}(B).
\end{equation}
The unused residual coordinates do not enter the second term. They do
enter the first term, and their bounds in $B$ cannot be omitted.

\subsection{The scale powers that pay for the estimate}

Put $H=\exp((\log\log s)^2)$ for the diverging monitor along the
comparison scales, and choose
\[
 \lambda=\frac1{4(nk+k+1)},\qquad
 B=H^\lambda,\qquad \xi=\varepsilon=H^{-2}.
\]
The parameter $s$ is taken sufficiently large that the logarithms are
defined and $H\ge1$. Since $\lambda nk\le1/4$,
\[
 C_{nk}(B)\le [2(nk+1)]^{nk}H^{1/4}.
\]
The right side of \eqref{eq:allarity-moment-error} is therefore at most
\[
 [2(nk+1)]^{nk}\{n(A+1)^n+A^n\}H^{-7/4}
                  \le H^{-3/2}
\]
for all sufficiently large $s$. The last threshold depends on the fixed
word and degree bound. It is uniform over the finite volume and replicas
once the preceding covariance and vertex hypotheses hold uniformly.

For completeness, the grade-tail purchase follows from the same powers.
If $r=H^b$, $a_i\le H^{b g_i+E}$, the linear-coordinate substitution
cost is at most $H^\lambda$, and $b\ge0$, $b\ge E+3$, then
\[
 a_i(H^\lambda)^k
   \le H^{b g_i+E+1/4}
   \le H^{-2}H^{b(g_i+1)}.
\]
The first inequality uses $\lambda k\le1/4$; the second follows from
$E+1/4\le b-2$. Thus every grade is treated at its own truncation
level while the arity and polynomial degree determine the common
Gaussian growth allowance.

\subsection{Taking the finite logarithm}

Use the squarefree algebra $\mathscr A_{\{1,\ldots,n\}}$ from
Section~\ref{sec:allarity-planted-compiler}. For a polynomial $q$ with
constant coefficient one, define
\[
 \log_*q=\sum_{j=1}^n\frac{(-1)^{j+1}}j(q-1)^j.
\]
Formal power-series identities truncated after degree $n$ give
$\exp_*(\log_*q)=q$ and $\log_*(\exp_*P)=P$ for $P$ with zero
constant coefficient. These identities can also be checked coefficientwise:
all powers above $n$ vanish, so the usual one-variable formal identities
in the nilpotent variable $q-1$ are finite identities here. In particular
the cumulant polynomial is uniquely determined by the moment polynomial.

We give an explicit local stability bound. Use the coefficient
$\ell^1$ norm on this algebra. If $|q-1|_1,|q'-1|_1\le M$, the
identity $u^j-v^j=\sum_{a=0}^{j-1}u^a(u-v)v^{j-1-a}$ yields
\[
 |\log_*q-\log_*q'|_1
 \le\left(\sum_{j=1}^n M^{j-1}\right)|q-q'|_1.
\]
If the coefficientwise moment error is at most $H^{-3/2}$, then
$|q-q'|_1\le2^nH^{-3/2}$. The free moment polynomial $q'$ varies
continuously on the compact comparison window, which is contained in
the distinct-point configuration domain, so its coefficient norm
is bounded. The moment error then bounds $|q-1|_1$ as well. Consequently
there is a constant $C$ for this fixed word and window such that
\begin{equation}\label{eq:allarity-cumulant-error}
 |\log_*q-\log_*q'|_1+H^{-3/2}
                  \le(C+1)H^{-3/2}\le H^{-31/32}
\end{equation}
for all sufficiently large $s$. The source uses the coefficient maximum
norm; the same conclusion follows from its being bounded by the
$\ell^1$ norm. The finite factor $2^n$ is absorbed only after $n$ is fixed.

On the actual intermediate Gaussian history, let $\mathcal S_{l,k}$
be the successive linear substitution and Gaussian integration operator
from level $k$ to level $l$. Define the moment polynomial by
\[
 q(z)=\sum_{I\subseteq\{1,\ldots,n\}}z_I\,
      \mathcal S_{l,k}\left(\prod_{i\in I}P_i\right)(y).
\]
Its constant coefficient is one, since every substitution and normalized
Gaussian integration preserves one. The connected polynomial obtained by
the finite partition formula is $\log_*q$: expanding the displayed
finite logarithm gives that partition formula directly, with coefficient
$(-1)^{j-1}(j-1)!$ for a partition into $j$ nonempty blocks.
Evaluation at the fixed intermediate field $y$ commutes with this finite
polynomial formula. This is the exact cumulant used in the planted
comparison. Therefore the moment estimate for the same reader
covariance and mean implies \eqref{eq:allarity-cumulant-error} for it.

If the planted local coefficient differs from this cumulant by at most
$H^{-3/2}$, its sum with the compensating negative free cumulant is
bounded by $H^{-31/32}$, by \eqref{eq:allarity-cumulant-error} and the
triangle inequality. Signed copies and complex conjugation preserve
these absolute coefficient bounds. This is the small initial remainder
needed for the subsequent marked propagation. Its use still requires
the local planted-coefficient error, the exact low-grade calibration,
and the common reader covariance estimates at the displayed powers;
they are not consequences of finite-dimensional logarithm continuity.
